\section{Exact two-label geometry and robust compatibility}\label{sec:binary47}
The positive Hankel normal form below retains the whole future array.
For two labels, a smaller exact normal form is available. It gives
both a class of finite certificates and an explicit error frontier.
The statements in this section concern the actual optimized numerical
error, not a sum of local defects or an average over a test law.

\subsection{Antisymmetric experiments}
We first allow any finite response table, not necessarily one obtained
from orthogonal matrices. Let the seeds be paired symbols
$X=\{+1,-1,\ldots,+n,-n\}$, let $w\in\calA^N$, and let
$j\in\{1,\ldots,d\}$ be the selected query. Suppose the target means
satisfy
\begin{equation}\label{eq:antisym47}
 f_{-i,w,j}=-f_{i,w,j},\qquad |f_{i,w,j}|\le1.
\end{equation}
The binary law is $(1+b f_{i,w,j})/2$. Write $\mathcal E_2(f)$ for the
minimum worst-word binary-TV error of a machine with two available
labels at \emph{every} cut $0,\ldots,N$. The final held answer also
has two labels. Random initialization and arbitrary cut-dependent
stochastic matrices are allowed. All extrema below are attained.

\begin{theorem}[An exact scalar normal form]\label{thm:scalar47}
For every table satisfying \eqref{eq:antisym47},
\begin{equation}\label{eq:scalar47}
 2\mathcal E_2(f)=
 \min_{\substack{|h_i|\le1,\ |\beta_{t,a}|\le1,\ |d_j|\le1}}
 \max_{i,w,j}
 \left|f_{i,w,j}-h_i d_j\prod_{t=1}^N\beta_{t,a_t}\right|.
\end{equation}
Thus the exact numerical optimization is the distance, in the uniform
norm, to the set of bounded decomposable tensors with one factor for
the positive seed, one for each command epoch, and one for the query.
Every feasible tensor on the right has a two-label stochastic
realization. For $N=0$ the empty product is one.
\end{theorem}
\begin{proof}
Represent a probability row on two labels by its signed coordinate
$r=2p-1\in[-1,1]$. Every stochastic update is an affine map
$r\mapsto\alpha+\beta r$ with $|\alpha|+|\beta|\le1$;
conversely this condition is sufficient. A bounded terminal mean has
form $c_j+d_jr$ with $|c_j|+|d_j|\le1$.
For a word $w$, the composite update has the form
$A_w+B_wr$, where $B_w=\prod_t\beta_{t,a_t}$.
If the two paired initialization coordinates are $r_i,r_{-i}$,
the odd part of the realized mean is therefore
\[
 \frac{\widehat f_{i,w,j}-\widehat f_{-i,w,j}}2
     =h_i B_w d_j,\qquad h_i=(r_i-r_{-i})/2.
\]
Taking the odd part cannot increase uniform error against the odd
target: each resulting residual is half the difference of two original
residuals. This operation is realizable without mixing whole machines.
Use initialization coordinates $\pm h_i$, replace every update by
$r\mapsto\beta_{t,a}r$, and replace the terminal column by $d_jr$.
All rows remain stochastic because $|\beta_{t,a}|,|d_j|\le1$.
This proves that the right side is no larger than twice the original
optimum, without retaining a symmetrization coin.

Conversely, initialize the positive label with probability
$(1\pm h_i)/2$, use the stochastic matrix
\begin{equation}\label{eq:binaryrow47}
 \frac12\begin{pmatrix}1+\beta_{t,a}&1-\beta_{t,a}\\
                       1-\beta_{t,a}&1+\beta_{t,a}\end{pmatrix},
\end{equation}
and terminal means $(d_j,-d_j)$. Induction gives the tensor in
\eqref{eq:scalar47}. Binary total variation is half the mean error.
Compactness of the factor cubes and continuity prove attainment.
\end{proof}

The scalar description of a two-state stochastic matrix is elementary.
The content needed here is that deleting all affine offsets produces
an admissible optimizer for an odd target, rather than merely a signed
linear approximation. No averaging over an unrecorded global random
choice is permitted. Tensor descriptions of stochastic processes are
classical; Ohta's reachable/null-space reductions concern a different
stationary observation model \cite{Ohta20}. Equation~\eqref{eq:scalar47}
keeps the prescribed command order, bounded query columns and fixed
numerical error objective simultaneously.

\subsection{Balanced sign certificates}
For an entry $e=(i,a_1,\ldots,a_N,j)$, form its incidence vector
$\iota(e)\in\mathbb F_2^M$, where
\[
 M=n+N|\calA|+d.
\]
There is one coordinate for each seed factor, each pair $(t,a)$, and
each query factor; $\iota(e)$ marks the factors used by that entry.
A set $C$ of nonzero target entries is \emph{balanced} if
$\sum_{e\in C}\iota(e)=0$. Equivalently each factor occurs an even
number of times. It is \emph{negative} if $\prod_{e\in C}f_e<0$.

\begin{theorem}[A finite sign obstruction and a gap dichotomy]\label{thm:sign47}
A negative balanced set $C$ certifies
\begin{equation}\label{eq:signlower47}
 \mathcal E_2(f)\ge\tfrac12\min_{e\in C}|f_e|.
\end{equation}
Suppose in addition that $f_e\in\{0,\rho,-\rho\}$ for a fixed
$0<\rho\le1$. Then exactly one of the following holds:
\begin{equation}\label{eq:dichotomy47}
 \begin{aligned}
 &\mathcal E_2(f)=\rho/2,
       &&\text{the nonzero signs do not factor;}\\
 &\mathcal E_2(f)\le\rho/4,
       &&\text{the nonzero signs factor.}
 \end{aligned}
\end{equation}
Factorability, and hence the alternative in \eqref{eq:dichotomy47},
is decided by a linear system over $\mathbb F_2$. In the first case
there is a negative balanced certificate of at most $M+1$ entries.
In the second case the proof constructs a legal two-label machine.
Complexity is polynomial in the \emph{explicit table size}; the table
can be exponential in the displayed horizon.
\end{theorem}
\begin{proof}
The product of any decomposable real tensor over a balanced set is
nonnegative, since each factor has even exponent. If every tensor
entry on $C$ approximated the target within a mean error strictly
smaller than $\min_C|f_e|$, it would be nonzero with the target sign.
Its product would be negative, a contradiction. Apply
Theorem~\ref{thm:scalar47} to obtain \eqref{eq:signlower47}.

Let $s_e\in\mathbb F_2$ encode the nonzero target sign. The signs
factor exactly when the system
\begin{equation}\label{eq:gf247}
 \ip{\iota(e)}{z}=s_e \pmod2\qquad(f_e\ne0)
\end{equation}
is consistent. If it is inconsistent, a linear combination of its
rows has zero left side and right side one. The corresponding entries
are negative and balanced. One may express a contradictory row in a
basis of at most $M$ earlier rows, so at most $M+1$ entries suffice.
Inequality~\eqref{eq:signlower47} gives $\mathcal E_2(f)\ge\rho/2$;
the fair-coin decoder attains the reverse inequality.

If \eqref{eq:gf247} has a solution, take the seed and epoch factors
to be the prescribed signs $\pm1$, and take each query factor to be
its prescribed sign times $\rho/2$. Every tensor entry then has
magnitude $\rho/2$ and all nonzero targets have the correct sign.
Both zero and nonzero entries have mean error at most $\rho/2$.
The construction \eqref{eq:binaryrow47} realizes these factors with
deterministic identity/swap command rows and bounded terminal means.
Thus $\mathcal E_2(f)\le\rho/4$. Ordinary finite-field elimination
both decides consistency and records the indicated certificate.
\end{proof}

The second alternative need not be an equality: a target which is
already a bounded decomposable tensor may have zero error. The theorem
classifies the strict-improvement threshold, not every smaller optimum.
The certificate is an elementary parity identity, not a new
nonnegative-rank theorem or a general polynomial-time minimization
algorithm for stochastic automata. The difference is useful: an
explicit witness certifies the error against \emph{every} two-label
machine, whereas an equivalence test checks a supplied machine.

\subsection{Two bottlenecks and an exact finite frontier}
Return to the orthogonal experiment. Put
\begin{equation}\label{eq:quarter47}
 X=\{\pm e_1,\pm e_2\},\qquad
 \calA=\{I,R\},\qquad
 R=\begin{pmatrix}0&-1\\1&0\end{pmatrix},
 \qquad0<\rho\le1/6.
\end{equation}
The finite signal range in \eqref{eq:quarter47} is separate from the
smaller sufficient calibration range of the arithmetic theorems.
Every target answer probability is in $[(1-\rho)/2,(1+\rho)/2]$.

\begin{theorem}[Two separated bottlenecks]\label{thm:bottleneck47}
At any horizon $N\ge1$, suppose a machine for \eqref{eq:quarter47}
has at most two labels at two distinct cuts $s<t$ in $\{0,\ldots,N\}$.
Its actual worst-word binary-TV error is at least $\rho/2$, regardless
of the widths at all other cuts. The fair-coin output attains this
bound at every such prescribed profile. The same conclusion holds
for any positive signal $\rho\le1$, independently of the upper
construction used below.
\end{theorem}
\begin{proof}
Freeze all commands before $s$ and after $t$ to $I$. Between the two
cuts retain the two words $I^{t-s}$ and $RI^{t-s-1}$, with $R$ in
the first executed position. Integrating the intermediate computation
gives two stochastic maps between the two endpoint alphabets. Pad a
one-label endpoint if necessary. The initial probabilities may depend
arbitrarily on the seed; the suffix gives one bounded column per query.
Thus the odd part of this restricted experiment has the one-step
scalar form of Theorem~\ref{thm:scalar47}, even when intermediate
alphabets are arbitrarily large.

Four entries of its positive-seed target are
\[
 f_{1,I,1}=\rho,\qquad f_{2,I,2}=\rho,\qquad
 f_{1,R,2}=\rho,\qquad f_{2,R,1}=-\rho.
\]
Each seed, intervening word and query occurs twice. These four entries
form a negative balanced set. Taking the odd part of any putative
approximation costs no extra error, so \eqref{eq:signlower47} gives
$\rho/2$. The zero mean decoder proves sharpness. No reset occurs
inside the interval and no command is replaced by a free input symbol:
collapsing the interval is a converse argument only.
\end{proof}

For $N=1$ let $\mathcal E(k_0,k_1)$ denote the optimum at the two
command-boundary widths. A positive minimum is attained even when the
allowed profile is infeasible for exact realization.

\begin{theorem}[A complete two-cut error frontier]\label{thm:frontier47}
In \eqref{eq:quarter47} at $N=1$, the exact frontier is
\begin{equation}\label{eq:frontiertable47}
\begin{array}{c|ccc}
 \mathcal E(k_0,k_1)&k_1=1&k_1=2&k_1\ge3\\\hline
 k_0=1&\rho/2&\rho/2&\rho/2\\
 k_0=2&\rho/2&\rho/2&\rho/4\\
 k_0\ge3&\rho/2&\rho/4&0
\end{array}
\end{equation}
Every upper bound has an explicit rational stochastic realization
when $\rho$ is rational. In particular, minimizing either cut
separately at tolerance $\rho/4$ permits two labels, but no single
machine can attain both minima simultaneously.
\end{theorem}
\begin{proof}
A one-label bottleneck erases all seed dependence; for the identity
word, opposite coordinate means force error at least $\rho/2$.
Theorem~\ref{thm:bottleneck47} treats $(2,2)$.
If either cut has two labels, then, for the identity word, the realized
four seed means lie on one affine line $n\cdot y=c$. If its uniform
coordinate error is at most $\delta$, then
\[
 \delta\|n\|_1\ge
 \max_{y\in\{\pm\rho e_1,\pm\rho e_2\}}|n\cdot y-c|
 =\rho\|n\|_\infty+|c|
 \ge(\rho/2)\|n\|_1.
\]
This proves the $\rho/4$ lower bounds after dividing mean error by two.
A degenerate segment lies in such a line as well.

Here are matching machines. For $a>0$, set
\[
 V_a=\{(a,0),(0,a),(-a,-a)\},\qquad
 \beta_a(y)=\frac13
 \left(1+\frac{2y_1-y_2}{a},\
       1+\frac{2y_2-y_1}{a},\
       1-\frac{y_1+y_2}{a}\right).
\]
Whenever $\beta_a(y)\ge0$, it sums to one and has barycenter $y$
in $V_a$.

For $(3,3)$ initialize with $\beta_{2\rho}(\rho x)$.
For each vertex $v\in V_{2\rho}$ and command $U\in\{I,R\}$,
use the row $\beta_{6\rho}(Uv)$ on $V_{6\rho}$.
The six rows are nonnegative: substituting the vertices gives
$2(Uv)_1-(Uv)_2\ge-6\rho$, the analogous second inequality,
and $(Uv)_1+(Uv)_2\le6\rho$.
Use the coordinates of $V_{6\rho}$ as terminal means. They lie in
$[-1,1]$ because $6\rho\le1$. The barycenter identity gives exactly
$\rho Ux$ for each word and seed.

For $(2,3)$ let the two initial labels represent
$\pm v$, where $v=(\rho/2,\rho/2)$; initialize $e_1,e_2$ at $+v$
and their negatives at $-v$. On command $U$, use the row
$\beta_{2\rho}(\pm Uv)$ and decode with $V_{2\rho}$.
All four rows are nonnegative. The output is
$U[\rho(x_1+x_2)(1,1)/2]$, differing from $\rho Ux$ in each
coordinate by at most $\rho/2$ for both $I$ and $R$.

For $(3,2)$ initialize with $\beta_{2\rho}(\rho x)$ and decode the
two terminal labels as $(1,1)$ and $(-1,-1)$.
For $v\in V_{2\rho}$ use positive-label probability
$(1+\ell(Uv))/2$, where $\ell(y)=(y_1+y_2)/2$.
Here $|\ell(Uv)|\le2\rho\le1$. The output mean is
$\rho((Ux)_1+(Ux)_2)(1,1)/2$, again with coordinate error $\rho/2$.
All remaining entries in the table follow by padding normalized unused
rows, or by the zero decoder. Every coefficient just displayed is
rational when $\rho$ is rational.
\end{proof}

For $\rho/4\le\varepsilon<\rho/2$, the feasible two-cut profiles
are exactly the upward closure of $(2,3)$ and $(3,2)$.
For $0\le\varepsilon<\rho/4$, both command cuts need at least three
labels; for $\varepsilon\ge\rho/2$, one at each command cut suffices.
The held answer is always charged separately as two labels. These are
finite numerical frontiers, not asymptotic claims about an irrational
alphabet. They complement rather than replace the growing-width laws.

\begin{theorem}[An arbitrary-horizon profile frontier]\label{thm:sparsefront47}
For the alphabet and seeds in \eqref{eq:quarter47}, let $N\ge1$ and
suppose every command-cut width belongs to $\{1,2\}\cup\{4,5,\ldots\}$.
For any $0<\rho\le1$, the exact optimized error is
\begin{equation}\label{eq:sparsefront47}
 \mathcal E(K_\bullet)=
 \begin{cases}
  \rho/2,&\text{some $K_t=1$, or at least two cuts have $K_t=2$},\\
  \rho/4,&\text{exactly one cut has $K_t=2$ and none has width one},\\
  0,&\text{every $K_t\ge4$}.
 \end{cases}
\end{equation}
The constructive machines have deterministic initialization and command
rows and rational decoder probabilities when $\rho$ is rational.
The excluded width-three profiles are not classified by this statement.
\end{theorem}
\begin{proof}
For a one-label cut freeze every word to the identity; opposite seeds
force error $\rho/2$. Two two-label cuts are treated by
Theorem~\ref{thm:bottleneck47}. With one such cut, freeze the entire
word to the identity. Its four output means lie in a segment, so the
line argument in Theorem~\ref{thm:frontier47} gives error $\rho/4$.

For the matching construction without a bottleneck, store which of
$\{\pm e_1,\pm e_2\}$ is the current vector. The command $R$ permutes
these four states, and the decoder is $\rho$ times the vector.
If the unique two-label cut is $s$, track this vector until that cut.
At cut $s$ replace the current vector $y$ by
$P y=(y_1+y_2)(1,1)/2$. There are only two possibilities,
$\pm(1,1)/2$, for these four axis vectors. This is a deterministic
compression, with no retained old phase. Thereafter four labels track
the quarter-rotation orbit of $Py$. The first transition out of the
bottleneck embeds its two possibilities into these four labels and
applies the actual command; subsequent transitions are permutations.
At $s=0$ perform the compression in initialization; at $s=N$ decode the
two possibilities directly. The decoder multiplies by $\rho$.
For each axis vector, $\|y-Py\|_\infty=1/2$; every subsequent quarter
rotation preserves this norm. The final mean error is exactly
$\rho/2$, proving the second line of \eqref{eq:sparsefront47}.
A fair-coin decoder handles the first line. Available excess labels
are padded by normalized unused rows. No random branch or clock
history is stored for free.
\end{proof}

\begin{corollary}[Stability under arbitrary target perturbations]\label{cor:stable47}
If a finite binary response table is perturbed by at most $\eta$ in
uniform row TV, its optimal error at any fixed profile changes by at
most $\eta$. In particular the two-bottleneck lower bound becomes
$\max(0,\rho/2-\eta)$ for any such perturbation of
\eqref{eq:quarter47}. The construction at $(2,3)$ or $(3,2)$ has
error at most $\rho/4+\eta$. Thus a strict compatibility gap persists
whenever $\eta<\rho/8$, including full-support perturbations.
\end{corollary}
\begin{proof}
For a fixed machine, the two target errors differ by at most $\eta$
by the triangle inequality; minimizing in both directions proves the
first claim. Apply it to \eqref{eq:frontiertable47}; the gap between
$\rho/2-\eta$ and $\rho/4+\eta$ is positive exactly when
$\eta<\rho/8$. Antisymmetry need not survive the perturbation.
\end{proof}
